\section{LTI-Systeme in der realen Welt}
\label{sec:irl_lti}
Die Eigenschaften von LTI-Systemen sind in der digitalen Audiobearbeitung besonders interessant, wo Realweltsysteme als LTI-System interpretiert
und durch ihre Impulsantwort approximiert werden können.
Ein bekannter Anwendungsfall ist der Faltungshall (Engl.: \textit{Convolution Reverb}), der das Nachhallverhalten eines Raumes durch Faltung einer in dem Raum
gemessenen Impulsantwort nachstellt. Die Akustik eines Raumes wird durch einen Schallimpuls angeregt, mit einem Mikrofon aufgenommen, und digital gespeichert. 
Aus dem Messergebnis wird die Impulsantwort des Raums berechnet. Wird diese im Anschluss mit einem beliebigen Eingangssignal gefaltet, dann approximiert das
Ergebnis die Anwendung der Transferfunktion des Raums auf das Eingangssignal.
Der Raum selbst weist lineare Eigenschaften auf, die eingesetzten Schallwandler und
Signalverstärker verursachen aber leichte nonlineare Verzerrungen. Die auf die Schallwandlung folgende Ausbreitung durch 
die Raumluft ist widerum nicht vollständig Zeit-Invariant \cite{farina2000,farina2007a,prawda2024}. 
Durch geeignete Messverfahren in Kombination mit entsprechenden Verfahren zur Berechnung der Impulsantwort (\textit{Entfaltung}, engl. \textit{Deconvolution}) können auch Systeme mit leicht
nonlinearen und Zeit-Varianten Eigenschaften vollständig charakterisiert werden \cite{farina2000,farina2007a}. 
